% Generated from paper/data; do not hand-edit numeric cells.
\begin{table*}[!tbp]
\centering
\caption{Dataset roles, class prompts, and prediction filtering.}
\label{tab:dataset_protocol}
\footnotesize
\setlength{\tabcolsep}{2pt}
\begin{tabularx}{\linewidth}{@{}L{0.14\linewidth}L{0.10\linewidth}*{2}{C}L{0.20\linewidth}C@{}}
\toprule
Dataset & Role & Classes & Prompts & Background & \shortstack{Prediction\\threshold}\\
\midrule
DLRSD & Source & 17 & 17 & None & Off\\
iSAID & Target & 15 & 15 & Ignored & Off\\
LoveDA & Dev. & 7 & 7 & background & 0.5\\
Potsdam & Target & 6 & 7 & clutter & 0.0\\
Vaihingen & Target & 6 & 7 & clutter & 0.0\\
UAVid & Target & 7 & 7 & background & 0.3\\
UDD5 & Target & 5 & 5 & background & 0.4\\
VDD & Target & 7 & 9 & background & 0.5\\
\bottomrule
\end{tabularx}
\par\vspace{3pt}
\begin{minipage}{\linewidth}
\footnotesize Source: DLRSD adaptation and source-domain evaluation. Dev.: LoveDA checkpoint selection. Target: independent target-domain evaluation. Every configured class participates in forward inference. Aliases expand the prompt count without changing the semantic class count. Off means filtering is inactive. The same prompts and dataset-specific threshold are used for all local methods and variants. Background fallback uses a separate threshold of 0.0.
\end{minipage}
\end{table*}
